\usepackage{geometry}
\usepackage[parfill]{parskip}
\usepackage{mathtools}
\usepackage{tcolorbox}
\usepackage{amsmath,amssymb,amsthm}              
\usepackage{graphicx}
\usepackage{enumitem}
\usepackage{tikz}
\usepackage{mathrsfs}  
\usetikzlibrary{calc}
\usepackage{tikz-cd}
\usepackage{xcolor}
\usepackage{booktabs}
\usepackage[linesnumbered,ruled,vlined]{algorithm2e} 
\usepackage{lineno}
\usepackage[colorlinks]{hyperref}
\usepackage{mdframed}
\usepackage[symbols,nogroupskip,sort=none]{glossaries-extra}
\renewcommand*{\glspostdescription}{\medskip}

\newtheorem*{theorem*}{Theorem}
\newtheorem*{proposition*}{Proposition}
\newtheorem*{corollary*}{Corollary}

\newcommand{\N}{\mathbb{N}}
\newcommand{\Z}{\mathbb{Z}}
\newcommand{\Q}{\mathbb{Q}}
\newcommand{\R}{\mathbb{R}}
\newcommand{\C}{\mathbb{C}}
\newcommand{\G}{\mathbb{G}}
\newcommand{\delayX}{\mathcal{X}}
\newcommand{\X}{\text{x}}
\newcommand{\W}{\text{w}}
\newcommand{\Rho}{\mathrm{P}}
\newcommand{\Lgr}{\mathcal{L}}
\mathchardef\Re="023C
\newcommand{\GLdelay}{\mathrm{GL}_{\mathrm{delay}}}    
\newcommand{\Hrose}{\mathscr{H}} 
\newcommand{\krose}{k_{\theta,\sigma}}


\theoremstyle{plain} 
\newtheorem{theorem}{Theorem}[section]
\newtheorem{lemma}{Lemma}[section]
\newtheorem{corollary}{Corollary}[section]
\newtheorem{proposition}{Proposition}[section]


\theoremstyle{definition} 
\newtheorem{definition}{Definition}[section]

\theoremstyle{remark} 
\newtheorem{remark}{Remark}[section]
\newtheorem{example}{Example}[section]
\newtheorem{assumption}{Assumption}[section]

\DeclareMathOperator{\diag}{diag}
\newcommand{\Var}{\mathrm{Var}}
\newcommand{\Cov}{\mathrm{Cov}}
\DeclareMathOperator{\E}{\mathbb{E}}

% a grey framed box with small padding
\newmdenv[
  linecolor=gray!60,
  linewidth=.8pt,
  topline=true,
  bottomline=true,
  skipabove=\baselineskip,
  skipbelow=\baselineskip,
  innerleftmargin=6pt,
  innerrightmargin=6pt,
  innertopmargin=4pt,
  innerbottommargin=4pt,
  backgroundcolor=gray!5,
  frametitlefont=\normalfont\bfseries,
]{assumptionbox}


% ------------------------------------------------------------------
% A numbered list tailored for assumptions:
%   (W1) ..., (W2) ...   or   (K1) ..., (K2) ...
% ------------------------------------------------------------------
\newlist{assumptionlist}{enumerate}{1}          % base = enumerate, depth = 1
\setlist[assumptionlist]{
    leftmargin=1.4cm,                          % indent
    labelsep=0.2cm,                            % space after label
    itemsep=0.4\baselineskip,                  % vertical gap between items
    label=\textbf{(\Alph*)}                    % default label (over-ridden per list)
}

\pagecolor{white}